\renewcommand{\bibname}{\DEEN{Quellenverzeichnis}{References}}
\bibliographystyle{plainnat} % I.4
\bibliography{references} % J.7

\ifIncludeLegalRegisters % J.8
  \addchap{\DEEN{Rechtsprechungsverzeichnis}{Table of Cases}}
  \TemplateHint{%
    Name des Gerichts (i.\,d.\,R. abgekürzt), Art der Entscheidung: Urt.
    (= Urteil), Beschl. (= Beschluss), G. (= Gutachten), vom (Datum der
    Entscheidung), Aktenzeichen, Fundstelle (Abkürzung, falls üblich),
    Erscheinungsjahr, ggf. Band, Teil, erste Seite bzw. Spalte des
    Entscheidungstextes.}
  % I.1-I.4: Eintraege alphabetisch in diese Umgebung setzen.
  \begin{hfhsourceentries}
    \hfhsourceentry{Name des Gerichts \dots}
  \end{hfhsourceentries}

  \addchap{\DEEN
    {Verzeichnis der Verwaltungsanweisungen und Parlamentaria}
    {Administrative and Parliamentary Sources}}
  \TemplateHint{%
    Bezeichnung der erlassenden Behörde/Stelle vom (Datum), Aktenzeichen
    (ggf. Betreff), Name der Fundstelle (übliche Abkürzung),
    Erscheinungsjahr der Fundstelle, ggf. Band, Teil, erste Seite bzw.
    Spalte.\par
    Institution-Drucksache, Drucksachennummer vom (Datum),
    Bezeichnung/Betreff der Drucksache, ggf. Seitenangaben.}
  % I.1-I.4: Eintraege alphabetisch in diese Umgebung setzen.
  \begin{hfhsourceentries}
    \hfhsourceentry{Behörde, Datum, Aktenzeichen, Fundstelle.}
    \hfhsourceentry{Institution-Drucksache, Nummer, Datum, Betreff.}
  \end{hfhsourceentries}
\fi

\ifIncludeAIRegister % J.9
  \addchap{\DEEN
    {Verzeichnis der eingesetzten KI-Werkzeuge}
    {Register of AI Tools Used}}
  \TemplateHint{\DEEN
    {Diesen Abschnitt nur behalten, wenn KI-Werkzeuge eingesetzt wurden. Die
     Einträge müssen den tatsächlichen Einsatz nachvollziehbar dokumentieren.}
    {Retain this section only if AI tools were used. Entries must document the
     actual use transparently.}}

  \begingroup
  \footnotesize
  \setstretch{1}
  \begin{tabularx}{\textwidth}{@{}YYYY@{}}
    \toprule
    \textbf{Name des verwendeten KI-Tools} &
    \textbf{Betroffene Teile der Arbeit} &
    \textbf{Einsatzform / Genutzte Funktion} &
    \textbf{Einsatztiefe bzw. Erläuterung, wie die Ergebnisse genutzt
      wurden} \\
    \midrule
    \textit{(Name des genutzten KI-basierten Hilfsmittels, ggf. mit
      Versionsangabe)} &
    \textit{(Auf welche Teile der eigenen Arbeit hat sich die Nutzung
      bezogen?)} &
    \textit{Funktion} &
    \textit{Kritische Prüfung und konkrete Verwendung} \\
    \bottomrule
  \end{tabularx}
  \endgroup

  \textit{Beispiele für Einsatzform / genutzte Funktion
  (nicht abschließend):}
  \begin{itemize}
    \item Vorschläge für Textgliederung
    \item Formulierungsvorschläge für einzelne Textpassagen
    \item Vorschläge für Übergänge zwischen Abschnitten
    \item Übersetzen von Fachartikel
    \item Wortwiederholungen vermeiden, Synonyme finden
    \item Beispiele finden
    \item Vorschläge für Visualisierungen
    \item Textkorrektur
    \item \dots
  \end{itemize}

  \textit{Beispiele für Einsatztiefe bzw. Erläuterung, wie die Ergebnisse
  genutzt wurden:}
  \begin{itemize}
    \item KI-Vorschlag diente lediglich zur Inspiration; finaler Titel wurde
      eigenständig entwickelt.
    \item Ergebnisse aus KI-gestütztem Brainstorming wurden selektiv
      übernommen und kritisch reflektiert.
    \item Vorschläge der KI dienten lediglich als Denkanstoß, tatsächlicher
      Text vollständig selbst formuliert.
    \item Vorschläge der KI anhand wissenschaftlicher Quellen überprüft und
      stark überarbeitet.
    \item Inhaltlich fehlerhafte KI-Antwort wurde erkannt und korrigiert.
    \item KI-Übersetzung wurde überprüft und mit Fachvokabular ergänzt.
    \item KI half beim Finden geeigneter Synonyme; Auswahl erfolgte unter
      Berücksichtigung des Kontexts manuell.
    \item KI-Hilfe beim Finden von Beispielen, aber Auswahl, Bewertung und
      sprachliche Einbettung eigenständig.
    \item KI-generierte Grafik wurde vollständig übernommen und als solche
      kenntlich gemacht.
  \end{itemize}
\fi

\ifIncludeAppendix % J.10
  \addchap{\DEEN{Anlagenverzeichnis}{List of Appendices}}
  \begin{tabularx}{\textwidth}{@{}lXr@{}}
    \textbf{\DEEN{Anlage 1}{Appendix 1}:} &
    \nameref{app:material-one}\dotfill & \pageref{app:material-one} \\
    \textbf{\DEEN{Anlage 2}{Appendix 2}:} &
    \nameref{app:material-two}\dotfill & \pageref{app:material-two}
  \end{tabularx}

  \appendix
  \chapter{\DEEN{Bezeichnung der Anlage}{Appendix Title}}
  \label{app:material-one}
  Fügen Sie hier nur Material ein, das für das Verständnis notwendig ist,
  aber den Lesefluss im Haupttext stören würde.

  \chapter{\DEEN{Bezeichnung der Anlage}{Appendix Title}}
  \label{app:material-two}
  Fügen Sie hier nur Material ein, das für das Verständnis notwendig ist,
  aber den Lesefluss im Haupttext stören würde.
\fi

% Diese Marke bestimmt die Gesamtzahl in der Fusszeile. Die abschliessende
% Erklaerung ist absichtlich ausgenommen und zeigt keine Seitenzahl.
\phantomsection
\label{LastNumberedPage}
\clearpage
